\documentclass[11pt,letterpaper]{article}

\usepackage[margin=1in]{geometry}
\usepackage[T1]{fontenc}
\usepackage{lmodern}
\usepackage{microtype}
\usepackage{enumitem}
\usepackage{amsmath}

\setcounter{page}{6}
\setlength{\parindent}{0pt}
\setlength{\parskip}{5pt}
\setlist{nosep,leftmargin=1.6em}

\begin{document}

\begin{center}
  {\Large\bfseries CRC Package Interface Specification}\\[3pt]
  {\large Preparation, model-matrix, initialization, E-step,
  capture-model M-step, outcome likelihood, outcome-model M-step, and
  public EM fitting interface}
\end{center}

\section{Preparation interface}

The internal function \texttt{prepare\_crc()} now owns the complete pre-fit
pipeline. It resolves the outcome distribution and control settings, validates
the public inputs, parses capture indicators, outcomes, formulas, and
misclassification specifications, performs cross-component compatibility
checks, constructs the model matrices, and initializes the state weights and
capture-cell counts. It returns an object of class
\texttt{crcfit\_preparation} containing the prepared model components and the
validated fitting settings.

\texttt{crc\_fit()} records the matched call and delegates preparation to
\texttt{prepare\_crc()}. It transfers the prepared components into its local
environment and calls \texttt{perform\_em\_crc()} to fit the model. The local
variables created through \texttt{list2env()} are
declared explicitly so that static R code analysis can identify their
bindings. This separation allows unit tests to exercise preparation and the
internal fitting components independently while the public function provides
the complete fitting path.

\section{Model-matrix infrastructure}

The package now converts parsed model specifications into complete capture-cell
grids and aligned observation-level states. The infrastructure includes the
capture and outcome design matrices, misclassification contributions, and a
direct mapping from each observation state to its capture cell. Construction
of these components is integrated into \texttt{prepare\_crc()}, which is
called by \texttt{crc\_fit()}. The capture- and outcome-model M-steps are
coordinated by the internal EM algorithm, which is now integrated into
\texttt{crc\_fit()}.

\subsection{Capture-model matrix}

\texttt{build\_capture\_matrix()} constructs the Cartesian product of all
capture histories, categorical levels used by the capture model, and latent
classes. For misclassified variables, true-category levels and their ordering
are obtained from the corresponding confusion matrices. The function returns:

\begin{itemize}
  \item \texttt{cells}, the complete capture-model cell grid;
  \item \texttt{matrix}, the design matrix generated from the parsed terms;
  \item \texttt{unobserved}, an indicator of all-zero capture histories.
\end{itemize}

With $J$ capture sources, $K$ latent classes, and categorical variables with
$C_1,\ldots,C_p$ levels, the grid contains
$2^J K\prod_{q=1}^p C_q$ rows.

\subsection{Initial observation states}

\texttt{build\_observation\_states()} repeats every observed unit once for
each latent class. It retains the variables needed by the capture model,
outcome model, and misclassification mechanisms. It also keeps aligned copies
of the standardized outcome and observed values of misclassified variables.
An observation identifier maps every expanded row back to its original unit.
Only variables used by the capture or outcome model are retained in the state
data. Recorded values of variables subject to misclassification are stored
separately for evaluation of their confusion-matrix contributions.

\subsection{Misclassified-variable expansion}

\texttt{expand\_misclass\_states()} expands the current states over every
possible true level of one misclassified variable. Each expanded row receives
the contribution
\[
  \Pr(\widetilde G_j=h\mid G_j=g)
\]
from its confusion matrix. When \texttt{true\_if} applies, the observed
category receives probability one and all other candidate categories receive
probability zero.

For several misclassified variables, separate confusion matrices are used and
their contributions are multiplied. This corresponds to conditional
independence of the misclassification mechanisms given their respective true
categories. Dependent mechanisms may be supported in a future extension.
These contributions are likelihood terms, not posterior probabilities, and
must not be normalized at this stage.

\texttt{expand\_all\_misclass\_states()} applies this expansion successively
to every specified variable. When no misclassification mechanism is supplied,
it assigns a contribution of one to every latent state.

\subsection{Outcome-model matrix}

\texttt{build\_outcome\_matrix()} applies the parsed outcome terms to the
fully expanded state data. Its rows therefore remain aligned with the state
data, standardized outcomes, observation identifiers, and accumulated
misclassification contributions. It returns \texttt{NULL} when no outcome
model is specified.

Every predictor in \texttt{outcome\_formula} is required to occur in
\texttt{capture\_formula}. This ensures that the predictor is represented in
the complete capture-cell grid and is consequently available for prediction
in the all-zero cells. The formulas may contain different terms; only the set
of variables is constrained.

\subsection{Capture-cell indices}

\texttt{build\_capture\_index()} returns one integer index for every expanded
observation state. Matching uses the capture indicators, categorical variables
in the capture model, and latent class. If the $i$th value is $r$, row $i$ of
the state data uses row $r$ of the capture-cell grid and capture design matrix.
The vector avoids repeated joins during EM iterations. Construction
fails if any state cannot be matched to the grid.

\subsection{Unified model-matrix object}

\texttt{build\_model\_matrices()} coordinates the complete construction
sequence and returns an object of class \texttt{crc\_model\_matrices} with:

\begin{itemize}
  \item \texttt{capture}, containing the cell grid, capture design matrix,
  and all-zero-cell indicator;
  \item \texttt{states}, containing expanded state data and aligned metadata;
  \item \texttt{outcome\_matrix}, containing the outcome design matrix or
  \texttt{NULL};
  \item \texttt{capture\_index}, linking state rows to capture-cell rows.
\end{itemize}

After parsing and compatibility validation, the preparation object stores this
component under the name \texttt{model\_matrices}. The value is retained by
\texttt{crc\_fit()} for interpretation and subsequent prediction.

\subsection{Cross-component validation}

Compatibility checks now reject outcome predictors absent from the capture
model. They also reject variables specified in \texttt{misclass} that are used
by neither the capture nor the outcome model. This prevents redundant state
expansion and makes the population-cell interpretation explicit.

\newpage

\section{Initialization}

\texttt{initialize\_crc()} creates normalized initial weights for the expanded
observation states and expected counts for the complete capture-cell grid. For
multiple latent classes, independent gamma draws produce a symmetric
Dirichlet allocation within each observation and candidate true-category
configuration. These allocations are multiplied by the misclassification
contributions and normalized within each observed unit. The concentration is
set by \texttt{control\$init\_alpha} and defaults to 20.

State weights are aggregated through the precomputed capture-cell indices.
Observed cells receive the corresponding weighted counts, while every
all-zero capture cell receives an initial expected count of one. The resulting
\texttt{crc\_initialization} object stores the state weights, cell counts, and
Dirichlet concentration. It is returned by \texttt{prepare\_crc()} and retained
by \texttt{crc\_fit()}.

\section{Control and progress interface}

\texttt{parse\_control()} resolves the user-supplied \texttt{control} list
against complete defaults and returns an object of class
\texttt{crc\_control}. The resolved object is stored by
\texttt{prepare\_crc()} and subsequently by \texttt{crc\_fit()}. Partial
overrides are supported, while unknown, duplicated, unnamed, and invalid
settings are rejected. The available settings are:

\begin{itemize}
  \item \texttt{init\_alpha}, the positive Dirichlet concentration, with
  default 20;
  \item \texttt{em\$max\_iter} and \texttt{em\$tolerance}, with defaults
  1000 and $10^{-6}$;
  \item \texttt{capture\$max\_iter} and
  \texttt{capture\$tolerance}, with defaults 100 and $10^{-8}$;
  \item \texttt{outcome\$max\_iter} and
  \texttt{outcome\$relative\_tolerance}, with defaults 1000 and $10^{-8}$.
\end{itemize}

The separate logical argument \texttt{verbose} defaults to \texttt{FALSE}.
It is validated and stored by \texttt{crc\_fit()} and accepted by the internal
EM driver. User-facing iteration messages remain to be implemented;
optimizer-specific tracing is not part of the public control interface.

\section{Expectation step}

\texttt{expectation\_crc()} implements the distribution-independent part of
the E-step. It accepts one fitted capture-model mean for every row of the
complete capture-cell grid and, when an outcome model is present, one outcome
log-likelihood contribution for every expanded observation state. A missing
outcome model is represented by zero log-likelihood contributions, equivalent
to likelihood contributions of one.

For state $(i,g,k)$, the unnormalized log weight is
\[
  \log m_{Y_i,g,k}
  + \log \Pr(\widetilde G_i\mid G_i=g)
  + \log L_V(i\mid g).
\]
Exact zero misclassification probabilities remain negative infinity on the
log scale and therefore produce posterior weights of exactly zero. The
function verifies that each observed unit has at least one admissible state
and normalizes the weights within unit after subtracting the largest log
weight. This log-sum-exp calculation avoids unnecessary underflow and
overflow.

The normalized state weights are aggregated through \texttt{capture\_index}
to update the expected counts for observed capture histories. Expected counts
for all-zero histories are set to their current fitted capture-model means.
The function returns a \texttt{crc\_estep} object containing one
posterior weight per expanded state and one expected count per capture cell.
It is called by \texttt{em\_iteration\_crc()} during every internal EM
iteration, including iterations initiated through the public fitting function.

\section{Capture-model maximization}

\texttt{maximize\_capture\_crc()} implements the capture component of the
M-step. Given expected capture-cell counts $n_c$ and capture-model rows $x_c$,
it maximizes
\[
  Q_Y(\theta)=\sum_c\{n_c\log m_c(\theta)-m_c(\theta)\},
  \qquad m_c(\theta)=\exp(x_c^\mathsf{T}\theta).
\]
The maximization is performed with \texttt{glm.fit()} using a Poisson family
and logarithmic link. Fractional expected cell counts are valid for this
objective because the omitted terms do not depend on $\theta$. The specific
Poisson warning about non-integer counts is suppressed, while other warnings
remain visible.

When no starting values are supplied, all coefficients start at zero, giving
initial fitted cell means of one. The function also accepts previous
coefficients as warm starts for later EM iterations. The public capture
controls are translated to \texttt{glm.control()} by
\texttt{em\_iteration\_crc()}.

The function rejects invalid counts or starting values, non-convergence,
rank-deficient design matrices, non-finite coefficients, and non-positive or
non-finite fitted means. It returns a \texttt{crc\_capture\_mstep} object with
the coefficients, fitted cell means, maximized objective value, iteration
count, and convergence indicator.

\section{Outcome log-likelihood}

\texttt{outcome\_loglik\_crc()} evaluates one outcome log-likelihood
contribution for every expanded observation state. Its arguments are the
aligned standardized outcome table, a positive mean \texttt{mu} for every
state, the outcome distribution, and, for the negative binomial model, a
single positive dispersion parameter \texttt{sigma}. Predictors in
\texttt{outcome\_formula}, including \texttt{.latent}, may therefore produce
different means for different candidate states of the same observed unit.

Let $W$ denote the corresponding untruncated count. Under
\texttt{"ztpois"}, $W\sim\operatorname{Poisson}(\mu)$. Under
\texttt{"ztnegbin"}, $W$ has mean $\mu$, size $1/\sigma$, and variance
$\mu+\sigma\mu^2$. In both cases the modeled outcome is the conditional
variable $V=W\mid W\geq1$. If $f$, $F$, and $S(q)=\Pr(W>q)$ denote the
untruncated mass, distribution, and survival functions, respectively, the
implemented contribution is
\[
  \ell_V=
  \begin{cases}
    \log f(v)-\log S(0), & \text{exact observation},\\
    \log\{S(l-1)-S(u)\}-\log S(0),
      & \text{interval }l\leq V\leq u,\\
    \log S(l-1)-\log S(0), & \text{right censoring }V\geq l,\\
    0, & \text{missing outcome}.
  \end{cases}
\]
Thus, interval bounds are inclusive and missing outcomes contribute a
likelihood factor of one under the ignorable coarsening assumption. Exact
contributions use the log probability mass directly. Interval contributions
use upper-tail log probabilities and \texttt{log\_sub\_crc()}, avoiding
unstable subtraction of nearly equal ordinary probabilities in the right
tail. No artificial probability floor is applied.

\subsection{Log-scale probability subtraction}

\texttt{log\_sub\_crc()} provides the stable operation
\[
  \log\{\exp(a)-\exp(b)\}, \qquad a\geq b,
\]
using \texttt{expm1()}. This helper is used for interval-censored outcome
likelihoods, where the required probability is obtained by subtracting two
survival probabilities. It also handles the case in which both probabilities
are zero and their log values are negative infinity.

\section{Outcome-model maximization}

The internal function \texttt{maximize\_outcome\_crc()} implements the outcome
component of the M-step by maximizing the posterior-weighted
expected log-likelihood
\[
  Q_V(\phi)=\sum_s \tau_s\ell_V(s;\phi),
\]
where $s$ indexes the expanded observation states, $\tau_s$ is the current
posterior state weight, and $\ell_V(s;\phi)$ is provided by
\texttt{outcome\_loglik\_crc()}. Missing outcomes have zero log-likelihood
contributions and therefore do not affect this maximization.

The function minimizes $-Q_V$ with \texttt{optim()} using
\texttt{method = "BFGS"}. For \texttt{"ztpois"}, the parameter vector
contains only the outcome regression coefficients $\beta$, with
$\mu_s=\exp(z_s^\mathsf{T}\beta)$. For \texttt{"ztnegbin"}, it also contains
\texttt{log\_sigma}; optimizing the logarithm guarantees
$\sigma>0$ without box constraints. No arbitrary coefficient bounds or
user-selected optimizer are exposed. The implementation currently uses
numerical derivatives.

Preliminary regression coefficients are obtained from a posterior-weighted
Poisson fit to exact outcomes when there are enough distinct exact
observations and the preliminary design matrix has full rank. Otherwise, the
intercept is initialized as the logarithm of the weighted mean of informative
lower bounds and the remaining coefficients are set to zero. For the negative
binomial model, a positive method-of-moments dispersion estimate based on
exact outcomes is used when at least two distinct exact observations are
available. The estimate is bounded below by $10^{-4}$, while 0.25 is the
conservative fallback. These quantities are starting values only and do not
replace maximization of the zero-truncated, censored objective.

The function accepts the preceding outcome coefficients and
\texttt{log\_sigma}, when applicable, as warm starts. Named starting vectors
are validated and reordered to match the outcome design matrix. The
\texttt{max\_iter} and \texttt{relative\_tolerance} elements of
\texttt{control\$outcome} are translated to the corresponding
\texttt{optim()} controls by \texttt{em\_iteration\_crc()}.

Only positive-weight states enter the weighted objective, which avoids
undefined products of zero weights and negative-infinite log likelihoods.
Invalid likelihood evaluations receive a finite penalty. An update is
accepted only when \texttt{optim()} reports convergence code zero, the
objective is finite, and $Q_V$ has not decreased beyond a small numerical
tolerance. A failed warm-start optimization is retried from the conservative
default values. The function also rejects rank-deficient informative design
matrices and non-finite or non-positive fitted means or dispersion values.

The returned \texttt{crc\_outcome\_mstep} object contains the regression
coefficients, fitted state means, dispersion for the negative binomial model,
maximized objective, optimizer evaluation counts, retry indicator, and
convergence indicator.

\section{EM algorithm}

\subsection{Single iteration}

\texttt{em\_iteration\_crc()} coordinates one complete EM update. It first
fits the capture model from the current expected capture-cell counts. When an
outcome model is present, it then fits that model using the current state
weights and evaluates the resulting outcome log-likelihood for every expanded
state. Finally, it calls \texttt{expectation\_crc()} with the fitted capture
means and outcome contributions. Models without an outcome skip the outcome
maximization and use zero log-likelihood contributions in the E-step.

The function translates the resolved capture and outcome controls to
\texttt{glm.control()} and \texttt{optim()} controls, respectively. It accepts
optional starting values for both M-steps. Its \texttt{crc\_em\_iteration}
result contains the \texttt{capture\_fit}, optional \texttt{outcome\_fit},
updated \texttt{state\_weights}, and updated \texttt{cell\_counts}.

\subsection{Iteration loop and warm starts}

The loop starts from the \texttt{crc\_initialization} object.
\texttt{perform\_em\_crc()} repeatedly applies the single-iteration update.
After the first iteration, the preceding capture coefficients are used as the
capture-model starting values. The outcome coefficients are likewise reused; for
\texttt{"ztnegbin"}, the preceding dispersion is appended as
\texttt{log\_sigma}. The loop stops at convergence or after
\texttt{control\$em\$max\_iter} iterations.

The returned \texttt{crc\_em\_fit} object contains the final component fits,
state weights, expected cell counts, number of iterations, convergence
indicator, and convergence diagnostics. Reaching the iteration limit returns
the final update with \texttt{converged = FALSE}.

\subsection{Public fitted object}

\texttt{crc\_fit()} stores the EM result as \texttt{em\_fit}, retains the
validated specifications, model matrices, initialization, call, and input
data, sets \texttt{fitted = TRUE}, and returns an object of class
\texttt{crc\_fit}. If the EM iteration limit is exhausted, the final result is
returned with its non-convergence indicator and a warning is issued. This
keeps the last estimates available for diagnostic inspection.

\subsection{Convergence criterion}

For iteration $t$, the loop calculates
\begin{align*}
  \delta_\tau^{(t)}
    &= \max_s |\tau_s^{(t)}-\tau_s^{(t-1)}|,\\
  \delta_Y^{(t)}
    &= \max_c |\log m_c^{(t)}-\log m_c^{(t-1)}|,\\
  \delta_V^{(t)}
    &= \max_s |\log \mu_s^{(t)}-\log \mu_s^{(t-1)}|,\\
  \delta_\sigma^{(t)}
    &= |\log \sigma^{(t)}-\log \sigma^{(t-1)}|.
\end{align*}
Here $s$ indexes expanded observation states, $c$ indexes capture cells,
$\tau$ denotes posterior state weights, $m$ denotes fitted capture means, and
$\mu$ denotes fitted outcome means. The outcome-mean term is omitted when no
outcome model is present, and the dispersion term is included only for
\texttt{"ztnegbin"}. Explicitly monitoring dispersion makes the package
criterion stricter than the simulation implementation, in which dispersion
was not a separate convergence component.

The overall change is the maximum of the applicable components. Convergence
requires this value to be finite and strictly smaller than
\texttt{control\$em\$tolerance}. The first iteration has infinite change
because no preceding fitted values exist, so at least two iterations are
required. Log differences provide scale-free, approximately relative changes
for positive model quantities; state probabilities are compared directly
because zero is admissible.

\section{Automated tests}

The public fitting interface, model-matrix, initialization, control, E-step,
both M-steps, outcome-likelihood, and EM components are covered by automated tests
implemented with the \texttt{tinytest} package. Separate test files cover
model matrices, misclassification expansion, cross-component compatibility,
initialization, control parsing, expectation calculations, maximization,
outcome log-likelihood calculations, EM orchestration, and public fitting.

The existing tests cover capture grids and design matrices, row alignment,
factor levels, all-zero cells, misclassification expansion, compatibility,
initialization, and E-step behavior. Test fixtures for these components now
call \texttt{prepare\_crc()} directly when testing parsed inputs, model
matrices, or initialization objects. The E-step and both M-steps remain tested
directly, while the dedicated EM tests exercise their internal orchestration.

Initialization tests verify normalized, reproducible state weights and their
capture-cell aggregation, including the single-class case and invalid control
values. Fifteen E-step assertions cover posterior normalization, impossible
states, cell aggregation, numerical stability, dimensions, and validation.

Thirteen control assertions cover defaults, partial nested overrides, the
\texttt{verbose} argument, unknown settings, malformed subsections, and
invalid values. Twelve capture-maximization assertions cover fitted-object
structure, dimensions, finite coefficients and means, objective evaluation,
warm starts, invalid inputs, and rank deficiency.

Fourteen outcome-maximization assertions cover the Poisson and negative
binomial models, coefficient names, positive fitted means and dispersion,
finite objectives, reordered named warm starts, retry from an invalid warm
start, invalid state-weight dimensions, and rank-deficient informative design
matrices. Together, the capture- and outcome-maximization test file contains
26 assertions.

The outcome-likelihood tests compare the Poisson and negative binomial results
with direct probability calculations for mixed exact, interval-censored,
right-censored, and missing outcomes. They also verify inclusive interval
bounds, a right-censoring boundary at one, numerical stability in the far
right tail, and validation of the mean and dispersion inputs. The complete
outcome-likelihood test file contains nine assertions.

The dedicated \texttt{test\_em.R} file contains 23 assertions. It verifies
single-iteration result classes, finite fitted values, positive negative
binomial dispersion, normalized state weights, and accepted warm starts. It
also exercises full-loop convergence for negative binomial and Poisson
outcomes, operation without an outcome model, the convergence-component
structure, and termination with \texttt{converged = FALSE} after a one-
iteration limit.

The dedicated \texttt{test\_crc\_fit.R} file adds 14 end-to-end assertions.
It covers a converged zero-truncated Poisson outcome model, a model without an
outcome component, the fitted object and internal component classes, retained
call and distribution information, and the warning and returned diagnostics
when the EM iteration limit is exhausted. The complete suite now contains 142
passing assertions.

\section{Planned next steps}

The core fitting path and its component-level and end-to-end tests are now in
place. Development should continue with the following main steps.

\begin{enumerate}
  \item Implement and test a \texttt{predict()} method for overall population,
  categorical-subpopulation, and latent-class estimates.
  \item Validate fitting and prediction using the accompanying IAOS submission
  data and simulation setting.
  \item Add user-facing methods such as \texttt{print()} and \texttt{summary()}.
  \item Extend the package documentation, \texttt{NEWS.md}, examples, and
  vignettes, and implement the remaining progress reporting.
\end{enumerate}

EM updates must follow the methodology in the accompanying paper and simulation
implementation. The existing automated tests establish the initialization,
matrix, alignment, iteration, and convergence invariants on which the public
fitting interface will depend.

\end{document}
